\section{Discussion}

\paragraph{Measured paths separate leverage from selectivity.}
Learned directions increase Target and Clean incidence, while their
neighbor- and capability-damage differences remain statistically unresolved.
A direction can therefore create more usable measured coefficients while its
collateral profile varies by layer. Layer 7 offers greater leverage together
with more collateral movement than layer 11. The measured clean regions show
where leverage and selectivity coincide: most contain one measured clean
coefficient, which motivates multi-strength evaluation alongside the primary
prediction labels.

\paragraph{Budget matching preserves the diagnostic hierarchy.}
Residual-norm matching aligns relative intervention magnitude, not response
rates: Qwen3.5 shows the largest gains and Ministral is intermediate. All three
models nevertheless reproduce more learned clean paths and a much larger
predictive contribution from low-dose response than static localization.


\paragraph{A small causal response is more actionable than static localization.}
Supervised geometry is valuable for constructing high-leverage directions and
describing their concentration and alignment, while these static descriptors
add little forecasting power by themselves. A strength-disjoint low-dose
response strongly predicts later outcomes across record, dataset, and domain
transfer. Localization becomes actionable when it is paired with a cheap causal
measurement of the specific path. PML therefore complements concept detection
and average steering evaluation
\citep{wu25a,silva-etal-2025-steering,goyal-iii-2026-steering} by testing
whether an internal signal supports effective and selective margin movement at
a chosen strength.

The weak, outcome-specific transfer of activation-path features to ROME defines
the scope of the result. A path can diagnose a fragile or promising activation
intervention, while persistent parameter editing remains a separate problem
\citep{NEURIPS2023_3927bbdc}.
